\documentclass{article}
\usepackage{graphicx} % Required for inserting images

\title{NY Aerial Imagery}
\author{Nicolas Abbate}
\date{April 2026}

\usepackage{etoolbox}
\AtBeginEnvironment{quote}{\par\singlespacing\small}
\usepackage{etoolbox}
\AtBeginEnvironment{quote}{\par\singlespacing\small}
\begin{document}
\section{Computer Vision and In-Batch Ordinal Pairwise Ranking Architecture}

At its core, the model takes as input an aerial image tile centered on a building footprint---together with a set of geographic covariates---and produces a single scalar score representing that building's relative structural wealth position at a given point in time. The full pipeline then aggregates these building-level scores to the census tract, maps their empirical quantiles to a target income distribution, and delivers zero-shot econometric estimates for any city without requiring additional labeled data.

The architecture is designed around a \textbf{dual objective}. First, it learns a monotonic ordinal continuum: the scalar score must strictly preserve the cross-sectional wealth ranking of buildings within a city at any given point in time. Second, it enforces temporal consistency: a physically stable building should occupy an invariant position on that continuum regardless of the year of observation, while a building that undergoes structural transformation should be displaced accordingly. Achieving both objectives simultaneously requires careful attention to which comparisons are econometrically valid.

\paragraph{Theoretical grounding: inverse spatial sorting.}
The model is motivated by the spatial sorting literature in urban economics \citep{ahlfeldt2015economics, diamond2016determinants}. In standard spatial equilibrium models, households sort across locations based on skill, intergenerational wealth endowments, and amenity preferences. Equilibrium prices capitalize the willingness-to-pay of the marginal bidder at each location, so the physical housing stock---its scale, quality, density, and morphology---reflects not merely architecture but the \textit{capitalized, permanent manifestation} of the long-run human wealth of its occupants. Human wealth here is understood in the permanent income sense: the expected present value of future earnings plus accumulated assets (skills and endowments).

This model inverts that mechanism. Rather than predicting sorting from fundamentals, it observes the durable physical capital of a neighborhood and recovers the baseline human wealth required to successfully outbid all others for that locational and structural bundle. Formally, the scalar output $\hat{R}_{b,i,t}$ can be interpreted as a \textit{latent monotonic mapping of human wealth}: it preserves the strict ordering of the present-discounted income expectations of the marginal occupant across buildings, without claiming to recover the cardinal dollar value of those expectations. This monotonic---but not cardinal---interpretation is precisely what the spatial equilibrium supports: the equilibrium pins down who wins the bid, not the exact magnitude of the winning bid.

\paragraph{The comparability constraint.}
This theoretical interpretation disciplines which comparisons the model may learn from. The network is a deterministic function of the image tile and its covariates. Given the same physical structure and the same local environment, it must produce the same score---this is not merely a desirable property but a logical consequence of the human wealth interpretation: a stable building in an unchanged neighborhood represents the same latent wealth bundle across time, so its mapping must be invariant. This gives rise to a strict comparability constraint on valid training pairs.

Two observations of the same building at different times---$(b, i, t)$ and $(b, i, t')$---are \textit{directly comparable}. If no structural change has occurred, the human wealth bundle is identical and the scores should agree; if change has occurred, the displacement in score space is identified by the change in the bundle itself. Two different buildings observed in the \textit{same} year---$(b, i, t)$ and $(c, j, t)$---are also \textit{directly comparable}, since both are positioned within the same city-wide income distribution at the same moment in time and the cross-sectional ordering is well-defined.

However, supervising the network with direct ACS label comparisons across \textit{both} dimensions simultaneously---$(b, i, t)$ against $(c, j, t')$ for $b \neq c$ and $t \neq t'$---is empirically unidentified. While Lemma 1 of our Representation Theorem guarantees that a true cross-spatiotemporal ordering exists in the latent wealth space, the observed ACS income distributions drift aggregately over time. A naive comparison of ACS labels across different buildings and different years conflates the true spatial wealth gradient with aggregate temporal drift. Therefore, the in-batch pairwise framework strictly masks out spatiotemporal pairs during training. The network is forced to \textit{learn} the cross-spatiotemporal bridge (Lemma 1) endogenously, by chaining together valid cross-sectional and temporal-trajectory gradients, rather than anchoring to historically drifted labels.

To enforce these objectives without introducing spatial cardinality mismatch or temporal label drift, we adopt an \textbf{In-Batch Pairwise Ranking} framework and abandon cardinal regression entirely. The loss supervises only the \textit{direction} of pairwise score differences, never their magnitude. This resolves the label-inconsistency problem inherent in longitudinal ACS data: under any magnitude-sensitive loss, a physically stable neighborhood artificially drifts downward in relative rank as the broader city's housing stock grows, producing corrupted gradients. By supervising direction alone, the model learns a transport function that connects how buildings behave over time without anchoring to historically corrupted magnitudes.

\subsection{Image Input and Spatial Spillover ($\tau$)}

Building footprint polygons are sourced from the official New York City Department of Information Technology and Telecommunications (DoITT) dataset. These polygons are used exclusively to identify building centroids; the polygon geometry itself is never fed to the network. Instead, for each building $b$, the model receives a raw multi-spectral aerial image tile centered on the building's centroid.

Because neighborhood wealth is shaped by localized externalities that extend beyond any single structure, the image tile is not clipped to the building footprint. We introduce a \textbf{spatial spillover hyperparameter, $\tau$}, defined as a buffer radius (in meters) around the centroid, so that the tile captures the immediate built environment surrounding the building. Let $X_{b,i,t}^\tau$ denote the aerial image tile of building $b$ in tract $i$ at year $t \in \{2008, 2010, 2012, 2014, 2016, 2018, 2020, 2022, 2024\}$, inclusive of the $\tau$-meter buffer.



\subsection{Feature Extraction: ScaleMAE with Low-Rank Adaptation}

The feature extractor is a \textbf{ScaleMAE} (Scale-aware Masked Autoencoder) backbone \citep{reed2023scalemae}, a Vision Transformer (ViT-Large/16) pretrained on multi-scale remote sensing imagery via masked image modeling. ScaleMAE is specifically designed to encode spatial resolution as a learned positional signal, making it well-suited for aerial imagery where the ground sampling distance (GSD) varies across acquisitions and zoom levels.

To adapt the pretrained backbone to the ordinal ranking task without catastrophic forgetting, we apply \textbf{Low-Rank Adaptation (LoRA)} \citep{hu2021lora} to the attention and MLP projection matrices of every transformer block. Specifically, trainable low-rank decomposition matrices of rank $r = 16$ are injected into the fused query-key-value attention projections (\texttt{qkv}) and the two MLP layers (\texttt{fc1}, \texttt{fc2}) of each block. The original pretrained weights remain frozen; only the LoRA adapters and the downstream regression head are updated during training. This reduces the trainable parameter count by over 99\% relative to full fine-tuning while preserving the rich spatial representations learned during pretraining.

The ScaleMAE backbone encodes each image tile $X_{b,i,t}^\tau$ into a sequence of patch token embeddings. Mean-pooling across the spatial patch dimension yields a single feature vector:
\begin{equation}
h_{b,i,t} = \text{MeanPool}\!\left( f_\theta\!\left(X_{b,i,t}^\tau\right) \right) \in \mathbb{R}^{1024}
\end{equation}
where $f_\theta$ denotes the LoRA-augmented ScaleMAE backbone.

\subsection{Late Fusion Score Head}

A \textbf{Late Fusion Head} projects the pooled image embedding---concatenated with an optional geographic covariate vector---onto a single scalar score. The architecture consists of:
\begin{enumerate}
    \item A linear compression layer that maps the 1024-dimensional image embedding to 64 dimensions, followed by GELU activation, layer normalization, and 40\% dropout.
    \item A final linear layer that maps the concatenation of the 64-dimensional compressed image features and the metadata vector $\mathbf{m}_{b,i,t}$ (currently a single covariate: distance-to-city-center) to the scalar output.
\end{enumerate}

The scalar output is:
\begin{equation}
\hat{R}_{b,i,t} = g_\phi\!\left(h_{b,i,t},\; \mathbf{m}_{b,i,t}\right)
\end{equation}
where $g_\phi$ denotes the late fusion head with weights $\phi$.

This scalar $\hat{R}_{b,i,t}$ represents the building's position on the ordinal poor-to-rich continuum. Within a census tract $i$, buildings share the same underlying socioeconomic context, so in expectation $\hat{R}_{b,i,t} = \hat{R}_{i,t}$ for all $b$. Deviations across buildings within the same tract are treated as measurement error. Because the model is never fed more than a few buildings per tract per epoch, it is regularized against overfitting to the visual idiosyncrasies of individual structures and is instead forced to learn the structural hierarchy of wealth at the neighborhood level.

\subsection{Out-of-Sample Generalizability and Quantile Mapping}

A core methodological commitment of this architecture is the intentional discarding of spatial cardinality at every stage of training and deployment. Urban morphologies scale differently: a one-standard-deviation increase in structural wealth in New York City manifests vertically, in building height and facade quality, whereas in Dallas the same gradient manifests horizontally, in lot size and setback. A model that learns absolute topological anchors from one city's training data will not transfer cleanly to another. By constraining the network to act as a pure monotonic ordinal ranker, we decouple the learned embedding from city-specific scaling.

Deployment to an out-of-sample city (e.g., Boston) proceeds as follows. The model computes scalar scores $\hat{R}$ for all target buildings based solely on their aerial imagery. Empirical quantiles of these scores are then mapped directly to the corresponding quantiles of the target city's ACS income distribution, delivering zero-shot econometric moment matching without any fine-tuning on labeled local data.

\subsection{In-Batch Pairwise Ranking: Batch Construction and Sampling Strategy}

The training pipeline constructs stochastically stratified \textbf{hybrid batches}, each designed to maximize pairwise supervision density from a single unified forward pass through the backbone.

\subsubsection{The Cross-Sectional Core ($\mathcal{B}_{cs}$)}
A \texttt{HybridBatchSampler} randomly selects a single anchor year $t$ and draws $B_{cs}$ distinct buildings observed in that year from the loaded cache shards. This cross-sectional core is the primary source of ordinal ranking supervision: every unique pair of buildings from different census tracts provides a valid directional constraint.

\subsubsection{The Temporal Auxiliary Set ($\mathcal{B}_{temp}$)}
For a fraction of the buildings in $\mathcal{B}_{cs}$ (controlled by the hyperparameter $\rho_{temp}$, default 0.20), the sampler queries the longitudinal panel for observations of the same building in a different year $t' \neq t$:
\begin{itemize}
    \item \textbf{Stable twins:} If building $k$ exists in $t'$ and underwent no structural change ($I_{k}^{valid}(t \to t') = 0$), the image for $t'$ is appended to the batch, contributing to the temporal stability penalty.
    \item \textbf{Changed twins:} If building $k$ exists in $t'$ but underwent statistically significant structural change ($I_{k}^{valid}(t \to t') = 1$), the image for $t'$ is appended to the batch, contributing to the temporal change ranking objective.
\end{itemize}

\noindent The complete hybrid batch is therefore $\mathcal{B} = \mathcal{B}_{cs} \cup \mathcal{B}_{temp}$, processed in a \textit{single forward pass} through the shared backbone.

\subsection{In-Batch Pairwise Ranking Loss}

A single forward pass encodes the entire hybrid batch $\mathcal{B}$ into a vector of scalar scores $\hat{R}$. The total loss is the decoupled sum of three objectives computed over valid subsets of these scores, plus a variance regularizer that prevents score collapse.

\subsubsection{Cross-Sectional Ranking ($\mathcal{L}_{cross}$)}

For the $B_{cs}$ cross-sectional images in year $t$, we enumerate all $\binom{B_{cs}}{2}$ unique, undirected pairs via the upper-triangular indices. Let $\mathcal{U}_{valid}$ denote the subset of unique pairs $(k, l)$ where buildings $k$ and $l$ reside in \textit{different} census tracts (same-tract pairs are masked out, as ACS labels provide no intra-tract variation).

Let $y_{k,l} = \text{sign}(Z_l - Z_k) \in \{-1, +1\}$ denote the true economic ordering from ACS labels, where $Z$ is the label (relative income score). We define the signed agreement as $\delta_{k,l} = y_{k,l} \cdot (\hat{R}_l - \hat{R}_k)$, which is positive when the model's ranking agrees with the ground truth.

The cross-sectional loss uses a \textbf{smooth logistic (RankNet) surrogate} rather than a hard hinge:
\begin{equation}
\mathcal{L}_{cross} = \frac{1}{|\mathcal{U}_{valid}|} \sum_{(k,l) \in \mathcal{U}_{valid}} \log\!\left(1 + \exp\!\left(-\frac{\delta_{k,l}}{T}\right)\right)
\end{equation}
where $T$ is a temperature parameter that controls the sharpness of the logistic sigmoid. Note that $\log(1 + e^x) = \text{softplus}(x)$.

\paragraph{Score-based temperature (preserving the ordinal framework).}
Critically, the temperature $T$ scales with the model's own output distribution, \textit{not} with ACS label magnitudes:
\begin{equation}
T = \max\!\left(m_{min},\;\; m_{base} \cdot \sigma_{batch}\right)
\end{equation}
where $\sigma_{batch} = \text{std}(\hat{R}_{cs})$ is the detached standard deviation of the cross-sectional scores within the current batch, $m_{base}$ is a scaling hyperparameter (default 1.0), and $m_{min}$ is a hard floor (default 0.05) that prevents division by near-zero values at initialization. Both $\sigma_{batch}$ and $m_{min}$ carry zero economic information---they are purely functions of the score head's output distribution. This design choice is essential to the ordinal framework: supervision is purely directional (sign of label differences), and the temperature merely adapts the gradient magnitude to the current prediction scale without injecting cardinal economic content.

\paragraph{Curriculum learning: easy-to-hard pair filtering.}
To improve gradient quality and prevent the model from fitting to noisy, ambiguous comparisons early in training, we implement an epoch-dependent curriculum that progressively admits harder pairwise comparisons. Each building is assigned to a score quintile bin $q \in \{1, 2, 3, 4, 5\}$ based on its ACS label. The absolute bin difference $|q_k - q_l|$ between paired buildings determines the difficulty of the comparison: large differences correspond to ``easy'' pairs with unambiguous ordinal direction, while small differences correspond to ``hard'' pairs where ACS measurement error may reverse the true ordering.

The curriculum schedule is:
\begin{itemize}
    \item \textbf{Epochs 0--59:} Only pairs with $|q_k - q_l| \geq 3$ (extreme quintile contrasts).
    \item \textbf{Epochs 60--149:} Pairs with $|q_k - q_l| \geq 2$ (adjacent quintile pairs admitted).
    \item \textbf{Epoch 150+:} Pairs with $|q_k - q_l| \geq 1$ (all inter-quintile pairs admitted).
\end{itemize}
Intra-quintile comparisons ($|q_k - q_l| = 0$) are \textit{permanently excluded}. The ACS margin of error makes within-quintile rankings pure noise; forcing a strict directional constraint on such pairs causes variance collapse as the model learns to shrink its output spread to trivially satisfy the loss.

\subsubsection{Temporal Stability Penalty ($\mathcal{L}_{stable}$)}

For buildings in the stable temporal subset $\mathcal{B}_{stable}$, the latent human wealth bundle is physically identical between years $t$ and $t'$. The model enforces temporal invariance via an unconditional L1 penalty. Each temporal twin is dynamically paired with its cross-sectional anchor by matching on building identifier (DOITT\_ID):
\begin{equation}
\mathcal{L}_{stable} = \frac{1}{|\mathcal{B}_{stable}|} \sum_{k \in \mathcal{B}_{stable}} \left| \hat{R}_{k,t} - \hat{R}_{k,t'} \right|
\end{equation}

\subsubsection{Temporal Change Ranking ($\mathcal{L}_{change}$)}

For buildings in the changed temporal subset $\mathcal{B}_{change}$, physical alteration identifies a shift in the human wealth bundle. Let $y_{k, t \to t'} \in \{-1, +1\}$ indicate the direction of the structural change (e.g., $+1$ if the tract gentrified). We supervise this displacement using the same smooth logistic surrogate used for cross-sectional ranking, applied to temporal pairs of the \textit{same} building:
\begin{equation}
\mathcal{L}_{change} = \frac{1}{|\mathcal{B}_{change}|} \sum_{k \in \mathcal{B}_{change}} \log\!\left(1 + \exp\!\left(-\frac{y_{k,t\to t'} \cdot (\hat{R}_{k,t'} - \hat{R}_{k,t})}{T_{change}}\right)\right)
\end{equation}
where $T_{change} = \max(m_{min}, \; m_{base} \cdot \sigma')$ and $\sigma'$ is the detached standard deviation computed over the union of cross-sectional and temporal change scores, ensuring the temperature reflects the full prediction-scale context.

\subsubsection{Variance Regularizer}

To prevent the score head from collapsing its output to a near-constant value---a failure mode in which the softplus loss can be trivially minimized by shrinking $\sigma_{batch}$ toward zero---we add a quadratic variance penalty computed exclusively over the cross-sectional scores:
\begin{equation}
\mathcal{L}_{var} = \left(\sigma_{cs} - 1\right)^2
\end{equation}
where $\sigma_{cs} = \text{std}(\hat{R}_{cs})$. This targets a unit standard deviation for the cross-sectional score distribution, ensuring the model maintains a well-spread prediction landscape. The penalty is applied only to cross-sectional scores (not temporal), maintaining internal consistency: both the temperature $T$ and the variance target reference the same $\sigma_{cs}$, eliminating any tug-of-war between temporal invariance and score spread.

\subsubsection{Total Objective}

The final unified objective aggregates the valid comparisons:
\begin{equation}
\mathcal{L}_{total} = \mathcal{L}_{cross} + \lambda_{s} \cdot \mathcal{L}_{stable} + \lambda_{c} \cdot \mathcal{L}_{change} + \mathcal{L}_{var}
\end{equation}
where $\lambda_s$ and $\lambda_c$ are weighting hyperparameters governing the relative influence of the temporal constraints against the cross-sectional ranking baseline.

\paragraph{Relation to the literature.} The cross-sectional ranking component of the loss is a smooth logistic pairwise surrogate, structurally identical to the RankNet formulation of \citet{burges2005learning}. The logistic surrogate is preferred over the hard hinge variant of \citet{joachims2002optimizing} because the smooth gradient propagates useful learning signal even for correctly ranked pairs, avoiding the dead-zone problem inherent in hinge losses where gradients vanish for pairs satisfying the margin constraint. The addition of an adaptive, score-based temperature $T$ that scales with the model's output distribution rather than with label magnitudes is, to our knowledge, novel in the satellite-imagery-for-economics literature. The combination of an unconditional L1 stability penalty on temporal positives with a directional logistic surrogate on negatives is structurally analogous to the Large Margin Nearest Neighbor objective of \citet{weinberger2009distance}, adapted from metric space to a 1-D scalar continuum. The curriculum learning schedule over quintile-based pair difficulty draws on the general curriculum learning framework of \citet{bengio2009curriculum}. To our knowledge, this specific combination---smooth ordinal supervision with score-based adaptive temperature, stable temporal anchors, directional structural change indicators, and progressive pair-difficulty curriculum---has not been applied previously in the satellite-imagery-for-economics literature, where existing work relies on cross-sectional regression \citep{jean2016combining} or cardinal pairwise differences \citep{yeh2020using} rather than a purely ordinal framework.

\subsection{Training Infrastructure}

\paragraph{Cyclic cache management.}
Because the full panel of aerial imagery for New York City exceeds available RAM, the training pipeline employs a \textbf{cyclic cache manager} that maintains a rotating window of pre-extracted, subsampled image shards on disk. Training proceeds over a fixed number of active shards; while the GPU processes the current shards, a background thread asynchronously generates the next batch of shards from the zarr-backed imagery store. This ensures the GPU is never starved for data while keeping RAM usage bounded.

\paragraph{Mixed-precision training and gradient accumulation.}
All forward passes are executed under \textbf{automatic mixed precision} (AMP) with float16 compute and float32 master weights, managed by a \texttt{GradScaler}. To achieve an effective batch size larger than the physical GPU memory allows, gradients are accumulated over $K$ micro-batches (default $K = 8$) before each optimizer step. This is equivalent to training with an effective batch of $B_{cs} \times K$ cross-sectional images, dramatically increasing the number of valid pairwise constraints per parameter update.

\paragraph{Optimizer and learning rate schedule.}
The optimizer is AdamW with weight decay 0.05. The learning rate follows a \textbf{ReduceLROnPlateau} schedule that monitors mean validation Spearman correlation, decaying by a factor of 0.8 after 20 epochs of no improvement, with a minimum learning rate floor of $\text{lr}/20$.

\paragraph{Data augmentation.}
Training tiles are extracted with a jitter pad of $\pm 10$ meters around the centroid, which is then realized as a random spatial crop at model resolution. Photometric augmentations include color jitter (brightness $\pm 0.3$, contrast $\pm 0.3$, saturation $\pm 0.2$), random grayscale (10\%), random horizontal and vertical flips. Validation and inference tiles receive no augmentation beyond ImageNet normalization.

\paragraph{Validation metric.}
Validation is performed using pointwise MSE loss for interpretable per-image tracking, complemented by \textbf{Spearman rank correlation} between predicted scores and ACS labels as the primary model selection criterion. The best model checkpoint is selected to maximize mean Spearman correlation across validation splits.

\subsection*{Summary of Core Architectural Hyperparameters}
\begin{itemize}
    \item \textbf{$\tau$ (Spillover Radius):} Buffer radius (in meters) around the building centroid defining the spatial extent of each aerial image tile. Automatically snapped to an exact subsampling multiple.
    \item \textbf{$\rho_{temp}$ (Temporal Fraction):} Fraction of the batch reserved for temporal auxiliary images (split between stable and changed twins). Default: 0.20.
    \item \textbf{$m_{base}$ (Margin Scale Factor):} Scaling coefficient for the adaptive temperature $T = \max(m_{min}, \; m_{base} \cdot \sigma_{batch})$. Default: 1.0.
    \item \textbf{$m_{min}$ (Margin Floor):} Hard minimum for the temperature, preventing numerical instability at initialization. Default: 0.05.
    \item \textbf{$\lambda_s$ (Stable Temporal Weight):} Weighting of the temporal stability L1 penalty. Default: 0.02.
    \item \textbf{$\lambda_c$ (Change Temporal Weight):} Weighting of the temporal change ranking loss. Default: 0.1.
\end{itemize}

\end{document}